\setlength\cwii{\b}
\setlength\cwiii{\c}
\setlength\cwi{\a}
\begin{tcbitemize}[raster columns=3,raster equal height=rows,
        raster force size=false,
        raster equal skip=-1.0pt,
        raster column 1/.style={width=\cwi,
            fontupper=\bfseries},
        raster column 2/.style={width=\cwii},
        raster column 3/.style={width=\cwiii,fontupper=\bfseries,
            halign=center,valign=center},
        ]
    %----------------------------------------------------------------
    \tcbitem[raster multirow=4] 1)
    \tcbitem[raster multicolumn=2,raster multirow=2,blankest]
        \begin{tcbitemize}[raster rows=2,raster columns=2,
               raster height=\tcbtextheight,
               raster column 1/.style={width=\cwii},
               raster column 2/.style={width=\cwiii,fontupper=\bfseries}
            ]
            \tcbitem Le noyau le plus stable est~: $\ce{^{37}_{17}Cl}$.
            \tcbitem (0,5pt)
            %--------------------------------------------------------
            \tcbitem \textbf{Justification~:} Le noyau $\ce{^{37}_{17}Cl}$
                possède l'énergie de liaison par nucléon la plus élevée.
            \tcbitem (0,5pt)
        \end{tcbitemize}
    %----------------------------------------------------------------
    \tcbitem 2.a)
    \tcbitem $\ce{^{36}_{17}Cl -> ^{36}_{18}Ar + ^{0}_{-1}e}$
    \tcbitem (1pt)
    %----------------------------------------------------------------
    \tcbitem 2.b)
    \tcbitem Désintégration $\beta^{-}$.
    \tcbitem (0,5pt)
    %----------------------------------------------------------------
    \tcbitem 2.c)
    \tcbitem
    \begin{align*}
        \Delta E&=\Bigl[m(\ce{^{36}_{18}Ar})+m_{e}-
        m(\ce{^{36}_{17}Cl})\Bigr] \times c^{2} \\
        &=\Bigl[35,967545+0,000549-35,968312\Bigr] \times 931,5=
          \qty{-0.203}{\MeV} \\
        \Rightarrow\; E_{\text{libérée}}&=|\Delta E|=\qty{0.203}{\MeV}\\
    \end{align*}
    \tcbitem (1pt)
\end{tcbitemize}

